\documentclass[a4paper]{article}
\usepackage[italian]{babel}
\usepackage[margin=20mm]{geometry}
\usepackage{array,tabularx,xcolor}
\usepackage[hidelinks]{hyperref}

\definecolor{fixdblue}{HTML}{1E3A52}
\newcolumntype{L}[1]{>{\raggedright\arraybackslash}p{#1}}
\newcolumntype{Y}{>{\raggedright\arraybackslash}X}
\setlength{\parindent}{0pt}
\setlength{\parskip}{5pt}
\renewcommand{\arraystretch}{1.3}

\begin{document}

\begin{center}
  {\LARGE\bfseries\color{fixdblue} Analisi SWOT di Fixd\par}
  \smallskip
  {\large Confronto con Google Maps e ProntoPro\par}
  \bigskip
  
\end{center}

Fixd permette di cercare servizi locali confrontando annunci, tariffe e
fasce di disponibilità. La consultazione è pubblica; per pubblicare o
contattare un professionista occorre essere autenticati.
Google Maps è un'alternativa per la ricerca locale; ProntoPro è un
concorrente diretto basato su richieste e preventivi.

\section*{Matrice SWOT}

Forze e debolezze sono interne al progetto; opportunità e minacce
riguardano il mercato e i fornitori. Le ipotesi restano da validare.

\medskip
{\small
\begin{tabularx}{\linewidth}{@{}L{1.5cm}YY@{}}
  \textbf{Origine} & \textbf{Fattori favorevoli} & \textbf{Fattori sfavorevoli} \\
  \hline
  \textbf{Interna} &
  \textbf{Forze progettuali, da validare (S)}\par
  \textbf{S1. Ricerca locale accessibile.} Ricerca pubblica con raggio in chilometri e filtro per categoria.\par
  \textbf{S2. Annunci con tariffa e calendario.} Informazioni utili per valutare il servizio prima del contatto;
  le fasce non costituiscono una prenotazione.\par
  \textbf{S3. Ricerca assistita facoltativa.} La descrizione del problema può aiutare a individuare il servizio adatto.
  &
  \textbf{Debolezze (W)}\par
  \textbf{W1. Offerta locale da costruire.} Servono annunci pertinenti per ogni zona e categoria.\par
  \textbf{W2. Segnali di fiducia limitati.} Mancano recensioni e profili verificati; l'account non certifica le competenze.\par
  \textbf{W3. Aggiornamento manuale.} Le fasce sono affidabili solo se aggiornate dai professionisti.\par
  \textbf{W4. Dipendenza da integrazioni esterne.} AI e cartografia dipendono da servizi esterni; i suggerimenti vanno validati.\par
  \textbf{W5. Sostenibilità da definire.} Mancano un modello di ricavo
  e una stima dei costi di gestione e acquisizione degli utenti.
  \\
  \hline
  \textbf{Esterna} &
  \textbf{Opportunità (O)}\par
  \textbf{O1. Bisogno locale da verificare.} Le fasce consultabili potrebbero aiutare chi ha vincoli di orario.\par
  \textbf{O2. Reti di professionisti sul territorio.} Associazioni e reti di artigiani possono favorire l'offerta iniziale.\par
  \textbf{O3. Strumenti digitali disponibili.} Calendari esterni potrebbero semplificare gli aggiornamenti, in una fase successiva.
  &
  \textbf{Minacce (T)}\par
  \textbf{T1. Google Maps presidia la ricerca locale.} Schede, recensioni e contatti lo rendono un sostituto credibile.\par
  \textbf{T2. ProntoPro presidia la scelta del professionista.} Preventivi e recensioni aiutano a scegliere per lavori da definire nel dettaglio.\par
  \textbf{T3. Condizioni dei fornitori esterni.} Costi e limiti di AI e cartografia possono incidere sulla continuità del servizio.
  \\
\end{tabularx}\par}

% Inizio seconda pagina: contenuto conservato ed escluso dalla compilazione.
% \newpage
% \section*{Strategie ricavate dalla matrice}
% 
% La dispensa distingue tra una risposta reattiva, che adatta le risorse
% alle opportunità esistenti, e una proattiva, che usa le capacità del
% progetto per costruire nuove opportunità.
% Gli incroci indicano come applicare questi due approcci a Fixd.
% 
% \begin{description}
%   \item[SO (S1, S2 + O1).] Provare un lancio locale per piccoli interventi
%   domestici, mettendo in evidenza raggio, tariffa e fasce consultabili.
%   È la risposta reattiva al possibile bisogno di trovare qualcuno
%   compatibile con i propri orari.
%   \item[WO (W1, W3 + O2, O3).] Coinvolgere una rete locale di artigiani
%   per costruire la prima offerta e favorire l'aggiornamento manuale
%   delle fasce. Valutare la sincronizzazione con calendari esterni solo
%   dopo il pilota. La collaborazione è la componente proattiva della
%   strategia, perché crea le condizioni per offrire disponibilità
%   consultabili in una stessa zona.
%   \item[ST (S2, S3 + T1, T2).] Mostrare tariffa e fasce dichiarate nel
%   percorso di consultazione e contatto autenticato. Aggiungere poi l'AI
%   per individuare annunci pertinenti già presenti nel catalogo.
%   La differenziazione dipende dalla qualità di questi dati e dalla
%   rapidità con cui il cliente arriva a un contatto utile.
%   \item[WT (W2, W4, W5 + T1, T2, T3).] Valutare, come sviluppi successivi,
%   la verifica dei profili e la segnalazione di annunci inattendibili.
%   Stimare costi ricorrenti e possibili ricavi prima di ampliare il pilota. Mantenere
%   utilizzabile la ricerca per testo, raggio e categoria anche quando
%   il fornitore AI è indisponibile.
% \end{description}
% 
% \textbf{Priorità per il lancio.} La prima versione utilizzabile (MVP)
% comprende account e accesso, annunci, ricerca per testo e raggio, filtro
% per categoria, contatto autenticato e gestione/consultazione delle fasce.
% Il calendario è essenziale per verificare S2. La ricerca AI resta una
% fase successiva, come indicato nel documento sulle priorità.
% 
% \textbf{Validazione proposta.} Provare il servizio in un comune con due
% categorie, per esempio idraulici ed elettricisti. Misurare:
% \begin{itemize}
%   \item ricerche con almeno un risultato pertinente sul totale delle ricerche;
%   \item calendari confermati o aggiornati negli ultimi sette giorni sul
%   totale dei calendari dei professionisti con annunci;
%   \item sessioni di ricerca con avvio del contatto sul totale delle
%   sessioni di ricerca, distinguendo gli abbandoni durante l'accesso.
% \end{itemize}
% Sono metriche proposte, non dati già disponibili. Servono una rilevazione
% degli eventi e una data di aggiornamento del calendario, oggi non descritte
% nel modello ER. Un avvio del contatto non dimostra che il lavoro sia stato
% assegnato. Il pilota servirà anche a verificare se l'obbligo di account
% scoraggia il contatto e se tariffa e disponibilità compensano tale passaggio.
% 
% 
% 
% 
% Fine seconda pagina commentata.

\end{document}
